\documentclass[11pt,a4paper]{article}
\usepackage[utf8]{inputenc}
\usepackage[vietnamese]{babel}
\usepackage{amsmath,amssymb}
\usepackage{geometry}
\geometry{a4paper, total={180mm,260mm}, left=15mm, right=15mm, top=18mm, bottom=18mm}
\usepackage{tikz}
\usetikzlibrary{shapes.geometric, arrows.meta, positioning, calc}
\usepackage{tcolorbox}
\tcbuselibrary{skins,breakable}
\usepackage{xcolor}
\usepackage{multicol}
\usepackage{array}
\usepackage{fontawesome}
\usepackage{fancyhdr}

\pagestyle{fancy}
\fancyhf{}
\renewcommand{\headrulewidth}{0.5pt}
\renewcommand{\footrulewidth}{0.5pt}
\fancyhead[L]{\small\bfseries\color{teal}\faBook\ SỔ TAY KIẾN THỨC TOÁN 3}
\fancyhead[R]{\small\bfseries\color{teal}Lớp Toán Thầy Lĩnh -- loptoanthaylinh.vn}
\fancyfoot[C]{\small\color{gray}-- Sổ tay kiến thức: Bài 56. Nhân số có bốn chữ số với số có một chữ số --}

\definecolor{primary}{HTML}{005F73}
\definecolor{secondary}{HTML}{0A9396}
\definecolor{accent}{HTML}{EE9B00}
\definecolor{danger}{HTML}{AE2012}
\definecolor{boxbg}{HTML}{F4F6F8}

\begin{document}

% --- BANNER TIÊU ĐỀ ---
\begin{tcolorbox}[
  enhanced,
  colback=white,
  colframe=primary,
  arc=3mm,
  boxrule=1.2pt,
  left=4mm, right=4mm, top=2.5mm, bottom=2.5mm,
  drop shadow=black!10!white
]
  \centering
  {\Large\bfseries\color{primary}\faBookmark\ BÀI 56: NHÂN SỐ CÓ 4 CHỮ SỐ VỚI SỐ CÓ 1 CHỮ SỐ}\\[1.5mm]
  {\small\itshape\color{black!75}Kĩ thuật đặt tính, thứ tự nhân từ phải sang trái, nhân không nhớ và có nhớ không quá 2 lần trên 1 trang A4}
\end{tcolorbox}

\vspace{-2mm}


\noindent
\begin{minipage}[t]{0.49\textwidth}\vspace{0pt}
\begin{tcolorbox}[
  enhanced,
  colback=blue!4!white,
  colframe=blue!75!black,
  arc=2mm, boxrule=0.8pt,
  left=2mm, right=2mm, top=1.8mm, bottom=1.8mm,
  title={\footnotesize\bfseries\color{white}\faTimes\ 1. Quy trình Đặt tính rồi tính}
]
{\footnotesize
\begin{itemize}\setlength{\itemsep}{1.5pt}\setlength{\leftskip}{-2.5mm}
  \item \textbf{Đặt tính:} Viết thừa số có 1 chữ số dưới thừa số có 4 chữ số, thẳng hàng đơn vị. Viết dấu $\times$, kẻ vạch ngang.
  \item \textbf{Thực hiện phép nhân:} Lấy thừa số thứ hai nhân lần lượt với từng chữ số của thừa số thứ nhất từ \textbf{phải sang trái}:
  $$\text{Đơn vị} \to \text{Chục} \to \text{Trăm} \to \text{Nghìn}$$
  \item \textbf{Quy tắc nhớ:} Lấy tích của hàng đó cộng thêm số nhớ từ hàng trước (nếu có), viết chữ số đơn vị, nhớ chữ số chục sang hàng tiếp theo.
\end{itemize}
}
\end{tcolorbox}
\end{minipage}
\hfill
\begin{minipage}[t]{0.49\textwidth}\vspace{0pt}
\begin{tcolorbox}[
  enhanced,
  colback=teal!4!white,
  colframe=teal!75!black,
  arc=2mm, boxrule=0.8pt,
  left=2mm, right=2mm, top=1.8mm, bottom=1.8mm,
  title={\footnotesize\bfseries\color{white}\faKey\ 2. Kĩ năng nhân nhẩm số tròn nghìn}
]
{\footnotesize
\begin{itemize}\setlength{\itemsep}{1.5pt}\setlength{\leftskip}{-2.5mm}
  \item \textbf{Quy tắc nhẩm tròn nghìn:}
  Lấy chữ số hàng nghìn nhân với thừa số thứ hai, rồi viết thêm ba chữ số $0$ vào tận cùng bên phải tích.
  \item \textbf{Ví dụ nhẩm:}
  $$2000 \times 4 = \mathbf{8000}$$
  (Nhẩm: $2\text{ nghìn} \times 4 = 8\text{ nghìn} = 8000$).
  $$3000 \times 3 = \mathbf{9000}$$
\end{itemize}
}
\end{tcolorbox}
\end{minipage}

\vspace{-2.5mm}
% --- KHỐI 3: SƠ ĐỒ TƯ DUY TIKZ ---
\noindent
\begin{tcolorbox}[
  enhanced,
  colback=orange!4!white,
  colframe=orange!80!black,
  arc=2mm, boxrule=0.8pt,
  left=2mm, right=2mm, top=1.5mm, bottom=1.5mm,
  title={\footnotesize\bfseries\color{white}\faSitemap\ SƠ ĐỒ TƯ DUY: TIẾN TRÌNH 4 BƯỚC NHÂN PHẢI QUA TRÁI}
]
\centering
\begin{tikzpicture}[
  box/.style={draw, rounded corners=2pt, align=center, font=\scriptsize\bfseries, inner sep=2.5pt, minimum height=0.75cm},
  arr/.style={-{Stealth[length=2mm]}, thick}
]
  \node[box, fill=purple!15, draw=purple!80!black, text width=2.6cm] (b1) at (0, 0) {BƯỚC 1: ĐƠN VỊ\\$d \times k$\\Viết ĐV, nhớ chục};
  \node[box, fill=teal!15, draw=teal!80!black, text width=2.6cm] (b2) at (3.5, 0) {BƯỚC 2: HÀNG CHỤC\\$c \times k + \text{nhớ}$\\Viết chục, nhớ trăm};
  \node[box, fill=blue!12, draw=blue!80!black, text width=2.6cm] (b3) at (7.0, 0) {BƯỚC 3: HÀNG TRĂM\\$b \times k + \text{nhớ}$\\Viết trăm, nhớ nghìn};
  \node[box, fill=orange!15, draw=orange!80!black, text width=2.6cm] (b4) at (10.5, 0) {BƯỚC 4: HÀNG NGHÌN\\$a \times k + \text{nhớ}$\\Ghi toàn bộ kết quả};

  \draw[arr, purple!80!black] (b1) -- (b2);
  \draw[arr, teal!80!black] (b2) -- (b3);
  \draw[arr, blue!80!black] (b3) -- (b4);
\end{tikzpicture}
\end{tcolorbox}

\vspace{-2.5mm}
% --- KHỐI 4: VÍ DỤ & MẸO TRÁNH BẪY ---
\noindent
\begin{minipage}[t]{0.49\textwidth}\vspace{0pt}
\begin{tcolorbox}[
  enhanced,
  colback=green!4!white,
  colframe=green!60!black,
  arc=2mm, boxrule=0.8pt,
  left=2mm, right=2mm, top=1.8mm, bottom=1.8mm,
  title={\footnotesize\bfseries\color{white}\faCheckCircle\ Ví dụ thực hành chuẩn mực}
]
{\footnotesize
\textbf{1. Phép nhân có nhớ:}\\
$\begin{array}{r@{}l}
  1418 & \\
  \times\quad 4 & \\
  \hline
  \mathbf{5672} &
\end{array}$
\qquad
\begin{minipage}[c]{0.5\linewidth}
-- $4 \times 8 = 32$, viết $2$ nhớ $3$.\\
-- $4 \times 1 = 4$, thêm $3$ là $7$, viết $7$.\\
-- $4 \times 4 = 16$, viết $6$ nhớ $1$.\\
-- $4 \times 1 = 4$, thêm $1$ là $5$, viết $5$.
\end{minipage}

\textbf{2. Bài toán thực tế:}\\
Mỗi xe tải chở được $1250\text{ kg}$ hàng. $3$ xe tải như thế chở được:\\
$1250 \times 3 = \mathbf{3750\text{ kg}}$ hàng.
}
\end{tcolorbox}
\end{minipage}
\hfill
\begin{minipage}[t]{0.49\textwidth}\vspace{0pt}
\begin{tcolorbox}[
  enhanced,
  colback=red!4!white,
  colframe=red!70!black,
  arc=2mm, boxrule=0.8pt,
  left=2mm, right=2mm, top=1.8mm, bottom=1.8mm,
  title={\footnotesize\bfseries\color{white}\faExclamationTriangle\ Mẹo tránh bẫy trắc nghiệm}
]
{\footnotesize
\begin{itemize}\setlength{\itemsep}{1.5pt}\setlength{\leftskip}{-2.5mm}
  \item \textbf{Bẫy cộng nhớ trước khi nhân:}\\
  Tuyệt đối tuân thủ: \textbf{Nhân trước, cộng nhớ sau}.\\
  Không lấy $(1 + 3) \times 4 = 16$, mà phải lấy $4 \times 1 = 4$ rồi $+ 3 = 7$.
  \item \textbf{Bẫy nhân với chữ số 0:}\\
  $0 \times k = 0$. Nếu có số nhớ từ hàng trước chuyển sang thì viết chính số nhớ đó vào kết quả.
\end{itemize}
}
\end{tcolorbox}
\end{minipage}


\end{document}
